\section{Principle V: Freedom}

Freedom is not at the beginning, but at the end. It is not at the root, but in the flowers and the fruits of human nature or, better said, human virtue. A man is freer in the proportion that he is better. This evolution is necessary. Our men believed they credited themselves with the prize of the effort by a Declaration of their famous rights, by displaying in their city halls and their schools, in their ministries and churches that that prize is acquired without effort. But to display everywhere that that everyone is born a millionaire, would everyone earn a shadow of a million?

Will you say that it is a right to freedom? The right to the million would not be more futile.

Freedom, a metaphysical principle, is one thing: freedoms are something else. The will of the people, the sum of individual wills, is one thing; the general will, the expression of the general interest of a nation, a race, a country is another thing, totally different.

The freedom of the crowd is called folly, that of a stupid person, stupidity, that of the crook, crime, that of the traitor, treason, and so forth. To say that freedoms limit each other is totally correct in the savage state or the state of anarchy: that simply means that the strong oppress or exploit the weak as long as the weak do not revolt against the strong and that potential victims did not find a sure process to snipe at and strafe their tormenters at will. Freedom conceived as regime or as principle, is a generally painful chaos.

Citizens who have reflected a little therefore do not let themselves be intimidated by the effects of words. They define freedom as a power or a force, a force or power that is worth what its subjects and its object are worth. The eternal ridicule of the 18th and 19th century consists in taking freedom for a rule or for a goal when it is only a means or the material.

Freedom lives in few men.

Whoever says real freedom says authority. Testamentary freedom creates the authority of the head of the family. Communal or provincial freedom creates the real power of the social authorities who live and reside in the place. Religious freedom recognizes the authority of spiritual laws and the internal hierarchy of a religion. Trade union and professional freedom devotes the authority of the disciplines and rules to the interior of corporations and professional organizations.

These are the observable facts.

If, however, in place of observing, one withdraws into himself to reflect, he realizes that the very nature of positive freedom—that which is not conceived by the relationship to an obstacle to surmount—freedom is the power.

Social issues are exercised in society and have as their point of application not marble, as the power of a sculptor, nor the machine, as the power of a mechanic, but by men engaged with us in social life.

What then is a freedom? A power.

Whoever can do nothing at all is not free at all. Whoever can do rather poorly is poorly free. Whoever can do infinitely is also infinitely free. One of the forms of power is wealth. Another of its forms is influence, it is physical force, it is intellectual and moral force. On what are these various powers exercised in various ways? On men. And that power, to whom does it belong? To men. When human freedom is found at its highest point and it met human objects to which to apply and impose itself, what name does it take? Authority. An authority is therefore only a freedom reaching its perfection.

Far from the idea that authority contradicts the idea of freedom, it is, on the contrary, its completion and complement. The freedom of a father of a family is an authority. The freedom of a religious confession is an authority. There are still the authorities as freedom of an association, the freedom of a commune, of a given province.

When one wishes that the State respects private initiative in the economic order, one fundamentally demands that it respects what Le Play named with a beautiful expression: the social authorities. All real freedoms, defined and practical, are of authorities. Freedom opposed to authority; freedom which consists in not being \emph{affected} by others, but also to not \emph{affect} them, that neutral freedom is not compatible with nature nor with the order of life.



\flrightit{Posted on 2012-11-11 by Cologero }

\begin{center}* * *\end{center}

\begin{footnotesize}\begin{sffamily}



\texttt{Cologero on 2012-11-11 at 19:05 said: }

Once again, the intent of this piece is a counter-argument to Rousseau's claim that ``men are born free". On the contrary, Maurras asserts that freedom is an achievement, that it comes with maturity. Ironically, Maurras is close to Fichte's claim that a man is free to the extent that he can submit to the laws of Reason. The relationship between freedom and authority brings to mind Evola's claim that only the king can be said to be totally free.


\hfill

\texttt{Avery on 2012-11-11 at 23:39 said: }

It appears that by restating freedom as a product of power, rather than something necessarily tied up with the kshatriya revolt, Maurras corrects Rousseau better than the Perennialists do, and is rewarded with a definition more meaningful both for historical societies and modern political purposes.

The Roman Senators were always ``free" to believe and do what they wanted, and ``freemen" or franklins had their power recognized even in the most Traditional society.


\end{sffamily}\end{footnotesize}
